\documentclass{report}
\usepackage{amsmath,amssymb}
\setlength{\parindent}{0mm}
\setlength{\parskip}{1em}
\begin{document}
\begin{center}
\rule{6in}{1pt} \
{\large Zichao, ZD Di \\
{\bf Optimization-Based Approach for Tomographic Inversion from Multiple Data Modalities}}

Mathematics and Computer Science Division \\ Argonne National Laboratory \\ 9700 S Cass Avenue \\ Lemont \\ IL \\ 60439
\\
{\tt wendydi@mcs.anl.gov}\\
Sven Leyffer\\
Stefan Wild\end{center}

Fluorescence tomographic reconstruction can be used to reveal the internal
elemental composition of a sample while transmission tomography can be used
to obtain the spatial distribution of the absorption coefficient inside the sample.
In this work, we integrate both modalities and formulate an optimization
approach to simultaneously reconstruct the composition and absorption effect in
the sample. Furthermore, we applied multigrid-based optimization
framework to speedup the performance. Improvements of accuracy and
convergence rate have shown for several
examples.


\end{document}
